% !TeX program = xelatex
% 编译方式：使用 XeLaTeX 编译（ctex 宏包需要 XeLaTeX 或 LuaLaTeX）
\documentclass[11pt,a4paper]{ctexart}

\usepackage[margin=2.4cm]{geometry}
\usepackage{amsmath,amssymb}
\usepackage{booktabs}
\usepackage{array}
\usepackage{graphicx}
\usepackage{xcolor}
\usepackage{enumitem}
\usepackage{listings}
\usepackage{tikz}
\usetikzlibrary{arrows.meta,positioning}
\usepackage[hidelinks]{hyperref}
\usepackage{caption}

\captionsetup{font=small,labelfont=bf}
\setlist{itemsep=2pt,topsep=3pt,parsep=0pt}

\definecolor{codebg}{RGB}{246,248,250}
\lstset{
  basicstyle=\ttfamily\small,
  backgroundcolor=\color{codebg},
  frame=single,
  rulecolor=\color{lightgray},
  breaklines=true,
  showstringspaces=false,
  columns=fullflexible,
  xleftmargin=6pt,xrightmargin=6pt,
  aboveskip=8pt,belowskip=6pt,
}

\title{\bfseries 信息检索系统设计与实现\\[2pt]\large （信息与知识获取 · 作业二 实验报告）}
\author{}
\date{2026 年 6 月}

\begin{document}
\maketitle

\begin{center}
\small
\begin{tabular}{>{\centering\arraybackslash}m{2.2cm}>{\centering\arraybackslash}m{2.8cm}>{\centering\arraybackslash}m{1.4cm}>{\raggedright\arraybackslash}m{5.6cm}}
\toprule
姓名 & 学号 & 角色 & 分工 \\
\midrule
曹忠昊 & 2023211964 & 组长 & 系统架构与异步爬虫、倒排索引构建、整体整合 \\
侯钧译 & 2023211728 & 组员 & 检索算法（BM25 / VSM / 混合检索 / 伪相关反馈）与评价 \\
徐舸山 & 2023211038 & 组员 & 多模态检索（Chinese-CLIP、VLM 精排）与实验报告 \\
\bottomrule
\end{tabular}
\end{center}

\vspace{0.4em}

\begin{abstract}
\noindent
本作业实现了一个面向真实中文新闻语料的信息检索系统。系统以自研异步爬虫采集了 802 篇权威媒体新闻并本地存储，
在此之上构建倒排索引，实现了 BM25、向量空间模型（VSM）与三路融合的混合检索，并加入伪相关反馈进行查询扩展。
在基础的文本检索之外，系统进一步扩展至跨媒体方向：以中文原生的 Chinese-CLIP 将图像与中文文本映射到同一向量空间，
支持以文搜图与以图搜文，并针对实测中发现的「模态间隙导致固定阈值失效」问题，设计了一套召回—精排的两段式流程，
由视觉大模型对候选图逐张裁定相关性。检索质量同时以与 BM25 解耦的自动基准和网页端的人工逐条打分两种方式评价。
除多模态增强需要可选的在线大模型外，检索主体全程在 CPU 上运行。
\end{abstract}

\section{任务背景与总体设计}

信息检索所要解决的核心问题，是在一个规模可观的文档集合中，依据用户以自然语言表达的需求，将最可能相关的文档排在前列。
本次作业的基本要求为：以开源爬虫采集不少于 100 篇文档并本地存储，建立倒排索引，实现向量空间检索模型的匹配算法，
支持自然语言查询，结果按相关度由高到低给出相关度、标题、匹配片段、URL 与日期，并尽量对结果准确率进行人工评价；
扩展要求则鼓励尝试多媒体检索与算法优化。

本系统在满足上述全部基本要求的同时，做了三处实质性增强。其一，检索模型不止于 VSM，而是将 BM25、
VSM 与语义召回经倒数排名融合（RRF）整合为混合检索，并提供伪相关反馈的查询扩展。其二，评价不限于一个主观打分页，
另实现了一套与排序算法解耦的自动基准，以避免「以 BM25 的打分证明 BM25 更优」这类循环论证。其三，也是投入最多的部分，
为跨媒体检索：本系统未止步于以文搜图的基本可用，而是分析了中文 CLIP 在实际数据上检索不准的根因，并以视觉大模型加以根治。
系统整体的数据流如图~\ref{fig:arch} 所示。

\begin{figure}[ht]
\centering
\begin{tikzpicture}[
  node distance=9mm and 12mm,
  box/.style={draw,rounded corners=2pt,align=center,inner sep=4pt,minimum height=8mm,font=\small},
  >={Stealth[length=2mm]}
]
\node[box] (crawl) {异步爬虫\\(去重/限流/配图)};
\node[box,right=of crawl] (corpus) {语料 + 配图\\(本地存储)};
\node[box,above right=of corpus,xshift=4mm] (text) {倒排索引\\BM25 / VSM / 混合};
\node[box,below right=of corpus,xshift=4mm] (clip) {Chinese-CLIP\\图文向量索引};
\node[box,right=22mm of corpus,yshift=0mm] (web) {Web 服务\\检索 / 多模态 / 评价};
\draw[->] (crawl) -- (corpus);
\draw[->] (corpus) -- (text);
\draw[->] (corpus) -- (clip);
\draw[->] (text) -- (web);
\draw[->] (clip) -- (web);
\end{tikzpicture}
\caption{系统总体数据流：爬虫采集的语料与配图分别进入文本索引与图文向量索引，统一由 Web 服务对外提供检索、多模态与评价功能。}
\label{fig:arch}
\end{figure}

% 如需插入系统界面截图，取消下面注释并放入图片文件：
% \begin{figure}[ht]\centering
% \includegraphics[width=0.9\textwidth]{screenshot_search.png}
% \caption{检索结果页：每条结果给出相关度、标题、匹配片段、来源与日期，并附逐条相关性打分按钮。}
% \end{figure}

\section{语料采集与数据真实性}

检索系统的质量首先取决于语料。课程明确要求数据须为开源爬虫真实抓取，因此本系统未使用任何模板生成或人工编造的数据，
并在代码层面删除了早期版本残留的数据生成器。最终入库 802 篇新闻，全部来自权威媒体，去重后覆盖 27 个子域名，
来源分布为：人民网 217、中国新闻网 178、新华网 139、央视网 130、中国经济网 48、科技日报 43、光明网 42、澎湃新闻 3。
每篇文档均完整保留原始的 URL、来源与日期，可逐条回溯至真实页面。

爬虫按工程化思路实现，关键设计如下。抓取基于 \texttt{aiohttp} 协程并发，并发上限可配置，吞吐较高而 CPU 占用低；
种子从多家媒体的财经、社会、健康、教育、科技等频道出发，以广度优先方式发现正文链接，以保证来源与主题的多样性。
正文以 \texttt{readability-lxml} 抽取，剥离导航、广告与评论。去重分两层：URL 级去重防止同一页面重复入库，
正文级则计算 SimHash 指纹，在海明距离阈值内判为近似重复，此步过滤掉了大量各站转载的同源稿件。
为减轻目标站点压力，轮换 User-Agent、在请求间加入随机间隔并对单域名限流。抓取逐篇落盘，同一 URL 覆盖同名文件，
故中途中断可续跑，重跑亦不产生冗余下载。正文中的 \texttt{<img>} 会被解析，按 PNG/JPG 文件头校验后下载，作为后续多模态检索的图像素材。

需说明的是，文档语料与以文搜图所用图库相互独立。新闻配图为从新闻页实际下载所得；
为提升以文搜图在灾害、科技等题材上的召回，另从图片搜索引擎补充了 325 张主题图，均为真实照片，并非编造或 AI 生成。
图库的扩充仅影响以文搜图，不改变文档语料的真实性。

\section{倒排索引与检索模型}

\subsection{倒排索引}

中文以 jieba 分词为基本单元，配合停用词过滤后建立倒排索引。索引记录每个词项到文档的倒排列表以及词频等统计量，
并以 pickle 持久化，启动时直接载入。当前词表规模为 47\,904 个词项。倒排索引是 BM25 与 VSM 共同的底层结构。

\subsection{BM25}

对查询 $Q=\{q_1,\dots,q_m\}$ 与文档 $D$，BM25 的打分为
\begin{equation}
\text{score}(D,Q)=\sum_{i=1}^{m}\text{IDF}(q_i)\cdot\frac{f(q_i,D)\,(k_1+1)}{f(q_i,D)+k_1\!\left(1-b+b\cdot\frac{|D|}{\text{avgdl}}\right)},
\end{equation}
其中 $\text{IDF}(q_i)=\ln\dfrac{N-n(q_i)+0.5}{n(q_i)+0.5}+1$，$N$ 为文档总数，$n(q_i)$ 为含 $q_i$ 的文档数，
参数取 $k_1=1.5$、$b=0.75$。相比朴素的 TF-IDF，BM25 一方面以 $k_1$ 控制词频饱和（词频再高，贡献亦趋于上限），
另一方面以 $b$ 进行文档长度归一化，避免长文档仅因词数多而获得偏高得分；这两点对新闻这类长度不均的语料尤为重要。

\subsection{向量空间模型}

VSM 将文档与查询均表示为 TF-IDF 加权向量，以余弦相似度衡量相关性：
\begin{equation}
\text{sim}(D,Q)=\frac{\vec{D}\cdot\vec{Q}}{\lVert\vec{D}\rVert\,\lVert\vec{Q}\rVert}.
\end{equation}
这是课程要求必须实现的匹配模型，也是与 BM25 对照的基线。

\subsection{混合检索与查询扩展}

实际检索中，纯词项匹配会受「词汇鸿沟」困扰：用户输入「汽车」而文档写作「轿车」，字面不匹配但语义相近。
为此并行执行 BM25、VSM 与一路基于 \texttt{sentence-transformers} 多语言向量的语义召回，再以倒数排名融合将三者合并：
\begin{equation}
\text{RRF}(d)=\sum_{r\in\{\text{bm25, vsm, sem}\}}\frac{1}{k+\text{rank}_r(d)},\qquad k=60 .
\end{equation}
RRF 只依据各路给出的排名而不直接混合分数，因而无需在不同量纲的得分之间做归一化，鲁棒性较好。三路召回互补：
BM25 偏重精确词项，VSM 偏重向量相似，语义召回补充同义与近义。

在此之上还实现了伪相关反馈（PRF）的查询扩展：将初次检索的前若干篇视为伪相关集，统计其中候选词的 $tf\cdot idf$ 权重，
剔除原查询词与过短词后，取权重最高的若干词加入查询并重检，以提升召回。前端设有开关控制是否启用，并回显扩展出的词，
便于观察效果。

\section{检索效果评价}

课程要求尽量对检索准确率进行人工评价，本系统在此基础上实现了自动与人工两套评价，互为补充。

\subsection{自动基准}

选取 10 个跨领域查询（人工智能、地震救援、新能源汽车、芯片制造、台风灾害、高等教育、医疗健康、自动驾驶、
自然灾害损失、财经投资），对 BM25 与 VSM 计算 P@5、P@10、F1@10、MAP、nDCG@10 与 MRR。
此处有一个易被忽视的陷阱：若相关性标签由待评测的排序算法自身计算得出，评测便沦为循环论证。
为规避此问题，相关性判定采用一套与 BM25 解耦的启发式——按正向关键词命中数划分相关等级（命中 $\ge 3$ 记部分相关，$\ge 5$ 记高相关），
出现跨领域负向关键词则降级，标题命中给予加成。10 个查询的平均结果见表~\ref{tab:bench}。

\begin{table}[ht]
\centering
\caption{BM25 与 VSM 在 10 个跨领域查询上的平均指标（自动基准）}
\label{tab:bench}
\begin{tabular}{lccc}
\toprule
指标 & BM25 & VSM & 优胜 \\
\midrule
P@5      & \textbf{0.5000} & 0.2800 & BM25 \\
P@10     & \textbf{0.3800} & 0.2100 & BM25 \\
F1@10    & \textbf{0.4529} & 0.2891 & BM25 \\
MAP      & \textbf{0.5518} & 0.4206 & BM25 \\
nDCG@10  & \textbf{0.5436} & 0.4307 & BM25 \\
MRR      & \textbf{0.7667} & 0.5202 & BM25 \\
\bottomrule
\end{tabular}
\end{table}

BM25 在全部指标上均优于纯 VSM，这与前文分析一致：词频饱和与长度归一化使其在新闻语料上更为稳健。
混合检索在 BM25 的基础上再融合语义召回，主要收益在于缓解词汇鸿沟。该基准可由 \texttt{python scripts/run\_evaluation.py} 复现。

\subsection{人工评价}

检索结果页的每条结果旁均设有「相关 / 部分相关 / 不相关」三个按钮。评分者点击后，系统实时计算本次查询的 P@5、P@10、MAP、
nDCG@10 与 MRR 并写入评价库，可在汇总页查看。这是面向真实使用者的主观评价，可弥补自动基准仅覆盖固定查询集的不足。

\section{跨媒体检索：中文 CLIP 与视觉大模型精排}

多模态是本系统投入最多的部分，其目标是支持用户以中文描述检索图片（以文搜图），或上传图片检索相关文档（以图搜文）。

\subsection{中文原生 CLIP 的选用}

最初采用多语言通用 CLIP（\texttt{xlm-roberta-base-ViT-B-32}），但其文本塔为多语言、视觉侧为英文预训练，
中文查询与图像的对齐属于折中方案。后改用中文原生的 Chinese-CLIP（\texttt{OFA-Sys/chinese-clip-vit-base-patch16}，
ViT-B/16 视觉塔加中文 RoBERTa 文本塔），其在约两亿中文图文对上训练，中文语义对齐明显更优，模型体积反而更小（约 600\,MB），
更适合纯 CPU 部署。系统保留多语言 CLIP 作为可切换的回退，便于对照。图文均被映射到同一 512 维空间，
以文搜图时以 FAISS 做余弦近邻召回，以图搜文时将上传图编码后与文档关联图比对。FAISS 索引以 \texttt{serialize\_index}
进行字节序列化持久化，绕开了 Windows 中文路径下 \texttt{write\_index} 失败的问题。

\subsection{图库的离线 VLM 质检}

从新闻页抓取的图像中混有大量「形为图像、实为噪声」的内容：栏目快讯的纯文字图、广告横幅、网站 logo、二维码、
网页或 App 界面截图、纯装饰素材等。仅凭「URL 关键词加最小边像素」的廉价规则无法识别，导致以文搜图的结果被大量无关图污染。
为此引入一道语义级质检：将每张图降采样至 512px，送入视觉大模型（通义千问 \texttt{qwen3-vl-plus}，经 DashScope 的 OpenAI 兼容接口调用），
由其判断是否为一张「有意义的新闻内容图」，仅输出 \texttt{\{keep, type, reason\}} 这一结构化结果，通过者方进入 CLIP 索引。

工程上对该步骤做了缓存与降级处理。判定结果按文件名持久化至 \texttt{data/vectors/vlm\_verdicts.json}，重复构建零调用、随交付一并分发，
因而即便批改环境无密钥，也能复现出「已清洗」的索引；多次失败时采取 fail-open（默认保留）且不写缓存，单次网络抖动不致误删图；
密钥仅从环境变量 \texttt{DASHSCOPE\_API\_KEY} 读取，不写入任何文件。最终的筛选流程为：846 张候选图先经廉价过滤剔除 149 张，
其余 697 张再送 VLM 质检剔除 178 张，留下 519 张内容图入索引。同一份判定结果还复用于文档详情页的配图展示与检索结果的图片角标计数，
以保证全站对「何为内容图」的口径一致。

\subsection{模态间隙与固定阈值的失效}

在图库清洗之后，以文搜图仍偶尔返回无关图。究其根源，在于 CLIP 的模态间隙：图向量与文向量在联合空间中分布于不同区域，
致使图文余弦被压缩在很窄的区间内（本库实测约 0.33--0.48）。更为关键的是，一些明显无关的图（机械零件、金字塔示意、纯文字标识）
的余弦可达 0.377--0.385，反而高于某些真正弱相关的图。这意味着无论将阈值设于何处，都会同时漏掉相关图、纳入无关图——
此为模型层面的固有现象，调整阈值仅能治标。

\subsection{召回—精排两段式流程}

既然单一阈值无法区分，便将检索拆为两段：CLIP 仅负责召回，先尽可能多地召回可能相关的图；相关性判定交由视觉大模型逐张精排。
召回阶段 CLIP 以一个宽松的下限（默认 0.28，仅用于滤除明显噪声、降低调用量）取 Top-N；精排阶段将每张候选图连同查询词送入 Qwen-VL，
由其从「以该词搜图的用户」视角判断是否相关，输出 \texttt{\{match, conf, desc\}}，仅保留 \texttt{match=true} 且置信度不低于 0.55 者，
再按置信度重排。判定按「查询词 + 文件名」缓存至 \texttt{data/vectors/vlm\_judge\_cache.json}，同一查询词再次检索时零调用、即时返回；
无密钥时退化为纯余弦阈值，保留可对照的旧行为。前端逐图显示模型的判定，例如「AI 判定相关 95\%（四川芦山地震震中示意图）」，
并统计「候选 N $\to$ 相关 M、剔除 K」，使大模型对结果的评价可见、可解释。开启精排后的实测见表~\ref{tab:judge}。

\begin{table}[ht]
\centering
\caption{开启 VLM 精排后的以文搜图结果（节选）}
\label{tab:judge}
\begin{tabular}{lccp{5.2cm}}
\toprule
查询 & CLIP 召回 & AI 判相关 & 典型保留结果 \\
\midrule
地震 & 12 & 12 & 芦山 7.0 级震中示意图、倒塌废墟搜救 \\
洪水 & 12 & 12 & 街道被淹乘舟救援、车辆涉水 \\
台风 & 12 & 12 & 巨浪拍岸、强风中树木弯折、卫星云图 \\
火灾 & 12 & 12 & 消防员扑救、山林野火 \\
新能源汽车 / 芯片 & 12 & 12 & 充电站车辆、晶圆芯片特写 \\
篮球 / 咖啡（库中无此题材） & 12 & 0 & 全部判为无关，如实返回空 \\
\bottomrule
\end{tabular}
\end{table}

离线质检（5.2）与在线精排互补：前者按「是否为内容图」清洗整库，后者按「与本次查询是否相关」对每次结果进行审定。
二者结合，使以文搜图的精度从依赖阈值的不确定判断，提升为逐张审定。

\section{对环境与社会可持续发展的影响}

可持续性是课程评分明确看重的方面，本文从两个角度加以考虑。在环境方面，异步并发与礼貌抓取减少了无效请求，
SimHash 去重避免了重复存储与重复索引，索引与向量均一次构建、反复复用，加之检索主体采用 CPU 友好的模型，
使系统能以较低的算力、带宽与能耗达到同等的检索效果。在社会方面，高效的检索提升了海量真实信息的可达性，
多模态检索进一步降低了普通用户获取信息的门槛。系统内设 \texttt{/sustainability} 页面对此作了集中说明。

\section{局限与展望}

语料虽已全部替换为权威媒体，但来源仍不够均衡，人民网约占 27\%、中国新闻网约占 22\%，后续可采用加权采样予以平衡；
灾害题材目前偏少，可针对应急管理、气象等频道定向补采。自动基准的相关性判定为启发式，更严谨的做法是积累人工标注、
构建标准测试集。在多模态方面，查询级精排依赖在线视觉大模型，虽已以缓存与无密钥降级缓解，但图库题材的覆盖仍可继续扩充；
对于垂直领域（如灾害图像），亦可对 CLIP 进行微调，以减少送审候选、进一步降低时延与调用量。

\section{运行与复现}

系统已附带构建好的索引与数据，开箱即可启动：

\begin{lstlisting}[language=bash]
pip install -r requirements.txt
python run.py web          # 启动 Web，默认 http://127.0.0.1:8090
# 从零构建（可选）：
python run.py crawl --max 800     # 异步爬取真实新闻
python run.py build               # 建立倒排索引
python run.py multimodal          # 构建 CLIP 图文向量索引
# 多模态 VLM 质检/精排为可选增强：设置环境变量 DASHSCOPE_API_KEY 后启用，
# 无密钥则自动降级为余弦阈值，并复用随附的已构建索引。
\end{lstlisting}

\noindent
启动后可访问检索、多模态、检索评价、数据看板与可持续等页面。自动基准以 \texttt{python scripts/run\_evaluation.py} 复现。

\end{document}
